\section{Methodology Charter}

\subsection{North Star}

The project will pursue decipherment only through constraints. We will not begin by assigning phonetic or semantic values to signs. We will begin by asking what the corpus permits.

\subsection{Core Rules}

\begin{enumerate}
  \item The cleaned corpus is the initial data layer. Claimed readings and translations are hypothesis material, not evidence.
  \item Any language-family claim must predict distributional behavior before it is used to explain individual signs.
  \item Iconographic resemblance is not a reading. It can generate a hypothesis, but not validate one.
  \item A proposed value must be tested against more inscriptions than the one that inspired it.
  \item Negative evidence matters: failed hypotheses stay in the ledger so we do not rediscover them under new names.
  \item We distinguish at least four levels: sign, allograph, sign cluster, and artifact formula.
  \item Archaeological context is part of the grammar. Site, medium, object type, iconography, completeness, and directionality are not metadata afterthoughts.
\end{enumerate}

\subsection{Data Tiers}

\begin{description}
  \item[Tier A] Complete inscriptions with known or strongly supported directionality, clean sign sequence, stable artifact metadata.
  \item[Tier B] Complete or near-complete inscriptions with uncertain directionality or minor damage.
  \item[Tier C] Damaged, fragmentary, unusual, or contextually uncertain inscriptions.
  \item[Tier D] Disputed, unprovenanced, translated-only, or otherwise circular records. These are excluded from baseline analysis.
\end{description}

\subsection{Initial Research Questions}

\begin{itemize}
  \item Which signs are strongly initial, medial, or terminal?
  \item Which signs behave like numerals, classifiers, titles, commodities, names, or formula delimiters?
  \item Do object types have distinct sign grammars?
  \item Does the system look more like language, restricted writing, emblematic identifiers, administrative coding, or a mixed system?
  \item Do any observed structures resemble Dravidian or Indo-Aryan morphology strongly enough to produce testable predictions?
\end{itemize}

\subsection{Decision Rule for Candidate Readings}

A reading is not allowed into the main model unless it satisfies all three conditions:

\begin{enumerate}
  \item It explains a distributional fact in the corpus.
  \item It predicts at least one additional distributional fact not used to invent it.
  \item It is more economical than credible non-linguistic and alternative linguistic explanations.
\end{enumerate}

\subsection{Bias Controls}

\begin{itemize}
  \item Maintain rival models in parallel.
  \item Score claims by evidence level.
  \item Prefer blind tests: train on one site or object class, test on another.
  \item Keep sign-number mappings distinct across Mahadevan, Parpola, Wells, and local datasets.
  \item Document every normalization step.
\end{itemize}

\subsection{Computational-First Protocol}

Because the project is led from a computer-science and data-science standpoint, claims should be expressed as scripts, features, scores, simulations, and reproducible outputs wherever possible. Linguistic ideas are useful when they become model constraints. The preferred workflow is:

\begin{enumerate}
  \item Convert an intuition into measurable variables.
  \item Generate corpus-wide features rather than isolated examples.
  \item Use simulations, resampling, or held-out tests when metadata permits.
  \item Treat model output as a ranking of hypotheses, not as proof.
  \item Promote a hypothesis only when the code, the corpus, and the visual/artifact record agree.
\end{enumerate}

Comparisons with Brahmic scripts are therefore allowed as hypothesis generators, especially for base-sign and modifier behavior, but not as assumed ancestry. A proposed Brahmic-style mechanism must predict a measurable pattern in the Indus corpus before it can influence interpretation.
